\begin{kbvrule}
\noindent{\kbvsecfont ವ್ಯಾಖ್ಯೆ :}\hspace{0.5em}ಪೂರ್ವಪದವು ಸಂಖ್ಯಾ ವಾಚಕವೂ ಉತ್ತರಪದವೂ ನಾಮಪದವೂ ಆಗಿರುವಾಗ, ಪೂರ್ವಪದವು ಉತ್ತರಪದದೊಂದಿಗೆ ಸೇರಿ ಆಗುವ ಸಮಾಸಕ್ಕೆ ದ್ವಿಗು ಸಮಾಸವೆಂದು ಹೆಸರು.
\end{kbvrule}

\begin{kbvderiv}
\vspace{4pt plus 3pt minus 1pt}
\begin{center}\begin{adjustbox}{max width=\textwidth}
{\renewcommand{\arraystretch}{1.22}\setlength{\tabcolsep}{4pt}
\begin{tabular}{l c l}
ಉದಾ : ನಾಲ್ಕು + ಮಡಿ & {\color{kbvgrey}$=$} & \kbvresult{ನಾಲ್ವಡಿ}{\color{kbvgrey} {[}ದ್ವಿ.ಸ.{]}} \\
ಎಂಟು + ದೆಸೆ & {\color{kbvgrey}$=$} & \kbvresult{ಎಣ್ದೆಸೆ}{\color{kbvgrey} {[}ದ್ವಿ.ಸ.{]}} \\
ಎರಡು + ಕೆಲ & {\color{kbvgrey}$=$} & \kbvresult{ಇರ್ಕೆಲ}{\color{kbvgrey} {[}ದ್ವಿ.ಸ.{]}} \\
ಒಂದು + ಕಟ್ಟು & {\color{kbvgrey}$=$} & \kbvresult{ಒಗ್ಗಟ್ಟು}{\color{kbvgrey} {[}ದ್ವಿ.ಸ.{]}} \\
ಒಂದು + ಕಣ್ಣು & {\color{kbvgrey}$=$} & \kbvresult{ಒಕ್ಕಣ್ಣು}{\color{kbvgrey} {[}ದ್ವಿ.ಸ.{]}} \\
\end{tabular}}
\end{adjustbox}\end{center}
\vspace{4pt plus 3pt minus 1pt}
\end{kbvderiv}

\vspace{4pt plus 4pt minus 1pt}\noindent{\kbvsecfont\fontsize{11pt}{14.5pt}\selectfont\color{kbvaccentdark} (೪)\hspace{0.45em}ಅಂಶಿ ಸಮಾಸ :}\hspace{0.5em}“ತುದಿಮೂಗು” ಎಂಬ ಸಮಸ್ತ ಪದವನ್ನು ವಿಗ್ರಹ ವಾಕ್ಯ ಮಾಡಿ\par

\indent ನೋಡೋಣ. ಮೂಗಿನ + ತುದಿ = ತುದಿಮೂಗು ಇಲ್ಲಿ ಪ್ರಧಾನವಾಗಿರುವುದು ಮೂಗಿನ ತುದಿಯೇ ಹೊರತು ಮೂಗಲ್ಲ “ತುದಿ” ಎಂಬುದು ಮೂಗಿನ ಒಂದು ಅಂಶ ಮಾತ್ರ. ಇಲ್ಲಿ ಮೂಗು ಅಂಶಿಯಾಗಿರುವುದು-\par

\indent - ಬ ಖಿ ಪ್ರ ವ ವಿಗ್ರಹ ವಾಕ್ಯದಲ್ಲಿ “ಮೂಗಿನ” ಎಂಬ ಪದ ಪೂರ್ವಪದವಾಗಿದ್ದರೂ, ಪೂರ್ವಪದ ಮತ್ತು ಉತ್ತರಪದಗಳು ಸೇರಿ ಸಮಾಸವಾದಾಗ ಸಮಸ್ತ ಪದದಲ್ಲಿ ಪೂರ್ವ ಮತ್ತು ಉತ್ತರಪದಗಳು ಅದಲು ಬದಲಾಗಿರುವುದನ್ನು ಕಾಣುತ್ತೇವೆ. ಆದುದರಿಂದ ಸಮಸ್ತಪದದಲ್ಲಿನ ಪೂರ್ವಪದವು ಪ್ರಧಾನವಾಗುವುದು. ಹೀಗೆ ಪೂರ್ವಪದದ ಅರ್ಥ ಪ್ರಧಾನವಾಗಿ ಉಳ್ಳ ಸಮಾಸಕ್ಕೆ ಅಂಶಿಸಮಾಸವೆಂದು ಹೆಸರು.\par

\begin{kbvrule}
\noindent{\kbvsecfont ವ್ಯಾಖ್ಯೆ :}\hspace{0.5em}ಪೂರ್ವ ಮತ್ತು ಉತ್ತರಪದಗಳು ಅಂಶಾಂಶಿಭಾವ ಸಂಬಂಧದಿಂದ ಸೇರಿ ಸಮಾಸವಾದಾಗ, ಪೂರ್ವಪದದ ಅರ್ಥವು ಪ್ರಧಾನ ವಾಗಿದ್ದರೂ ಅದು ಅಂಶಿ ಸಮಾಸವೆನಿಸುವುದು.
\end{kbvrule}

\indent ಅಂಶ--ಪೂರ್ಣದ ಒಂದು ಭಾಗ ; ಅಂಶಿ-ಅಂಶಗಳನ್ನು ಒಳಗೊಂಡಿರುವುದು.\par

\noindent\hspace*{2.2em}ಅಂಶಾಂಶಿಭಾವ--ಅಂಶಕ್ಕೂ ಅಂಶಿಗೂ ಇರುವ ಸಂಬಂಧ.\par
